\documentclass[11pt]{article}
\usepackage[margin=0.9in]{geometry}
\usepackage{amsmath,amssymb,booktabs,longtable,array,xcolor}
\usepackage[colorlinks,linkcolor=blue!45!black,urlcolor=blue!45!black]{hyperref}

\newcommand{\M}{\mathcal M}
\newcommand{\T}{\mathsf{T}}
\definecolor{good}{RGB}{20,110,60}
\definecolor{bad}{RGB}{160,30,25}
\definecolor{neut}{RGB}{95,95,110}
\definecolor{corr}{RGB}{150,95,10}

\newcommand{\POS}{\textcolor{good}{\textbf{Positive}}}
\newcommand{\NEG}{\textcolor{bad}{\textbf{Negative}}}
\newcommand{\NEU}{\textcolor{neut}{\textbf{Neutral}}}
\newcommand{\COR}{\textcolor{corr}{\textbf{Correction}}}
\newcommand{\FIN}{\textcolor{good}{Final}}
\newcommand{\RUN}{\textcolor{corr}{Still on the cluster}}

\newenvironment{todoitem}[2]{%
  \par\medskip\noindent\rule{\textwidth}{0.7pt}\par\nobreak\smallskip
  \noindent\textbf{\large TODO #1}\quad\textnormal{\small(draft #2)}\par\nobreak\smallskip
}{\par\medskip}

\newcommand{\ask}[1]{\noindent\textbf{Asked.}\ #1\par\smallskip}
\newcommand{\resp}[1]{\noindent\textbf{Response.}\ #1\par\smallskip}
\newcommand{\interp}[1]{\noindent\textbf{Interpretation.}\ #1\par\smallskip}
\newcommand{\verdict}[2]{\noindent\textbf{Sign.}\ #1\hfill\textbf{Status.}\ #2\par}

\setlength{\parindent}{0pt}
\setlength{\parskip}{4pt}

\title{Responses to the red TODO items}
\author{}
\date{27 August 2026}

\begin{document}
\maketitle

\section*{Summary}

Fifteen red items, which collapse to eleven distinct requests; four are the same
question asked once in the main text and once in the appendix. Nine are answered from data
already in hand, two require a change to a claim currently in the draft, and two numbers quoted
in the draft cannot be traced to any run and should be removed.

\begin{center}
\small
\begin{longtable}{@{}clll@{}}
\toprule
\# & Item & Sign & Status\\
\midrule
\endhead
1 & Three-panel figure, publication quality & \COR & \FIN\\
2 & Output-variety metric, integrator & \POS & \FIN\\
3 & ``Off-trajectory departure $0.26$'' & \NEG & \FIN\\
4 & Meaning of ``$4\times$ training'' & \NEU & \FIN\\
5 & Guidance formulation, weights, $\lambda$ & \POS & \FIN\\
6 & OpenCLIP / CIFAR-10 protocol & \NEU & \FIN\\
7 & The 27 empirical thresholds & \POS & \FIN\\
8 & ``$4\times$ training'' (appendix duplicate of 4) & \NEU & \FIN\\
9 & Capacity sweep at a second learning rate & \POS & \FIN\\
10 & Truncation sweep, $\varepsilon$ values and integral & \COR & \FIN\\
11 & Auxiliary-source perturbation sweep & \POS & \FIN\\
12 & Output variety (appendix duplicate of 2) & \POS & \FIN\\
13 & Guidance (appendix duplicate of 5) & \POS & \FIN\\
14 & OpenCLIP (appendix duplicate of 6) & \NEU & \FIN\\
15 & Super-resolution experiment & \POS & \RUN\\
\bottomrule
\end{longtable}
\end{center}

Two items need a decision from you rather than a number: item~3 and the
``factorization discrepancy $0.62$'' inside item~6/14. Item~10 obliges a rewrite of the
mechanism paragraph. Everything else can be pasted in.

% ------------------------------------------------------------------
\begin{todoitem}{1}{lines 38--41, Figure 1}
\ask{The publication-quality version of the three-panel coupling-ledger figure, together with
the underlying data.}

\resp{Attached as \texttt{fig1\_main.pdf} (vector), with \texttt{answers.json} carrying every
number behind it. Panels: (a) the flow-versus-do-nothing crossing at the operating point,
(b) the gap against width at each learning rate, (c) the guidance sign reversal against source
informativeness.}

\interp{The figure that existed in the notebook was \emph{not} the figure the caption describes.
The saved panel was a map of $\beta^\star(A)$ over six values of $j$: no crossing, no baseline,
no measured points. Panel~(a) as captioned did not exist and has now been built from
\texttt{beta\_star.csv} at $A=0.37$, $j=0.15$, plotting the analytic population landing law and
the analytic do-nothing level against the two measured curves. They cross at
$\beta^\star=0.6589$ (analytic) against $0.6778$ (measured). All output is now written as PDF as
well as PNG; the previous figures were $150$~dpi raster only.}

\verdict{\COR}{\FIN}
\end{todoitem}

% ------------------------------------------------------------------
\begin{todoitem}{2 and 12}{lines 47--48 and 265--268}
\ask{The definition of the output-variety metric, its values over the complete
integration-step sweep, the tested step counts, the integration scheme, and confirmation that
the one-shot regression baseline uses the same evaluation set and target-alignment definition.}

\resp{Variety is the \textbf{mean off-diagonal pairwise cosine among the six samples drawn for a
single condition}. Note the direction: the statistic \emph{decreases} as variety increases.
The solver is intrinsic Heun, second order, on a uniform grid, default $K=32$, swept over
$K\in\{1,2,4,8,16,32,64\}$. The one-shot regression baseline is scored by the identical routine
on the identical task, $768$ conditions $\times\,6$ samples --- confirmed.

\begin{center}
\small
\begin{tabular}{@{}rrrr@{}}
\toprule
$K$ & target cosine $\uparrow$ & pairwise cosine $\downarrow$ & polar $W_1$ $\downarrow$\\
\midrule
$1$  & $0.3541$ & $0.5299$ & $0.4623$\\
$2$  & $0.3524$ & $0.3897$ & $0.3286$\\
$4$  & $0.3389$ & $0.2593$ & $0.1632$\\
$8$  & $0.3254$ & $0.2030$ & $0.0611$\\
$16$ & $0.3197$ & $0.1858$ & $0.0360$\\
$32$ & $0.3175$ & $0.1838$ & $0.0327$\\
$64$ & $0.3172$ & $0.1829$ & $0.0321$\\
\bottomrule
\end{tabular}
\end{center}}

\interp{This is the cleanest demonstration of the metric inversion in the whole study. A single
integration step attains the \emph{highest} target cosine and the \emph{worst} distributional
fit, by a factor of $14$ in $W_1$. Choosing the integration budget by single-target cosine alone
selects the coarsest possible rollout. The effect should be attributed to Blau and Michaeli's
perception--distortion tradeoff, with our contribution being the closed-form quantification, not
the phenomenon.}

\verdict{\POS}{\FIN}
\end{todoitem}

% ------------------------------------------------------------------
\begin{todoitem}{3}{lines 81--82}
\ask{How the ``off-trajectory departure of $0.26$'' is defined and computed.}

\resp{No quantity by that name exists. A search of the sweep notebook, both retest notebooks, the
cluster script and the experiments source returns nothing that computes a trajectory departure,
and no logged register entry carries the value in that role.}

\interp{The nearest quantities in the same table are the source-conditioned in-plane energy
$0.367$ and its pairwise cosine $0.6559$; the analytic do-nothing ladder level is $0.26374$,
which is numerically closest to $0.26$ but is the \emph{source baseline}, not a departure, and
already appears elsewhere in the paper under its own name. My recommendation is to delete the
sentence rather than reverse-engineer a definition that fits the number. If you would prefer to
keep a statement of that kind, the defensible version is the in-plane energy: azimuth
conservation is exact for the population field ($1.000$), broken by training to $0.858$ for
$\M_1$ and to $0.367$ for $\M_2$.}

\verdict{\NEG}{\FIN}
\end{todoitem}

% ------------------------------------------------------------------
\begin{todoitem}{4 and 8}{lines 91--93 and 185--187}
\ask{What ``$4\times$ training'' means, and confirmation that all remaining training and
evaluation settings are unchanged.}

\resp{It is optimisation steps: $3{,}000\rightarrow12{,}000$, exactly four times. Everything else
is held fixed --- batch $256$, AdamW at $3\times10^{-4}$ under maximal-update parameterisation,
gradient clipping at $1.0$, the same checkpoint-selection criterion, the same evaluation
protocol and the same evaluation seeds. Confirmed by inspection of the training entry point:
the long run differs from the short one in the step count argument alone.}

\interp{Bookkeeping. It makes the comparison legitimate, which is the point of the question.}

\verdict{\NEU}{\FIN}
\end{todoitem}

% ------------------------------------------------------------------
\begin{todoitem}{5 and 13}{lines 102--103 and 289--293}
\ask{The exact classifier-free-guidance formulation, the guidance weights, the definition and
tested values of the source-informativeness parameter $\lambda$, the underlying terminal-cosine
values, and the unguided reference against which the reported changes are measured.}

\resp{The guided velocity is
$v_w=v_{\text{uncond}}+w\,(v_{\text{cond}}-v_{\text{uncond}})$, tangent-projected to the sphere,
with the unconditional branch obtained from the same network under the condition-dropout flag and
training carried out at $10\%$ condition dropout. \textbf{$w=1$ is the unguided reference}: the
implementation returns the conditional field unchanged at that value, so every reported change is
measured against it. Weights tested: $w\in\{1,2,3\}$.

$\lambda$ interpolates the initial state from pure noise toward the informative source,
$z_0\leftarrow\mathrm{slerp}(u,\,z_0,\,\lambda)$ with $u$ uniform on the sphere, so $\lambda=0$ is
a noise start and $\lambda=1$ a fully informative one; tested at
$\lambda\in\{0,0.25,0.5,0.75,1\}$, two seeds per cell. Change in terminal target cosine relative
to $w=1$:

\begin{center}
\small
\begin{tabular}{@{}rrrr@{}}
\toprule
$\lambda$ & $w=1$ & $w=2$ & $w=3$\\
\midrule
$0.00$ & $+0.0001\pm0.0012$ & $+0.0424\pm0.0002$ & $+0.0770\pm0.0004$\\
$0.25$ & $-0.0000\pm0.0012$ & $-0.0325\pm0.0003$ & $-0.0673\pm0.0007$\\
$0.50$ & $+0.0001\pm0.0014$ & $-0.1271\pm0.0003$ & $-0.2244\pm0.0001$\\
$0.75$ & $+0.0002\pm0.0014$ & $-0.2430\pm0.0138$ & $-0.3346\pm0.0065$\\
$1.00$ & $+0.0007\pm0.0011$ & $-0.3721\pm0.0173$ & $-0.4905\pm0.0039$\\
\bottomrule
\end{tabular}
\end{center}}

\interp{Two things worth putting in the text. First, the $w=1$ column is flat at zero to within
$0.0014$, which is the internal consistency check and should be labelled as the reference rather
than plotted as data. Second, and stronger than the current claim: the sign flips at
$\lambda_0=0.133$, not near $\lambda=1$. Guidance calibrated for a noise start begins to
\emph{hurt} as soon as the initialisation carries even slight target information. That is a
sharper and more useful statement than ``it reverses for informative sources''.

Note also that the draft quotes $+0.157$ and $-0.505$; the current run gives $+0.0770$ and
$-0.4905$ with seed error bars. The draft numbers are from an earlier pass. The newer values are
smaller on the positive side but carry standard errors of $4\times10^{-4}$, which is the more
defensible pair.}

\verdict{\POS}{\FIN}
\end{todoitem}

% ------------------------------------------------------------------
\begin{todoitem}{6 and 14}{lines 121--123 and 308--311}
\ask{The OpenCLIP architecture and checkpoint, the CIFAR-10 split and number of evaluated
images, the preprocessing and blur construction, the embedding normalisation, the definition of
the target-alignment metric, the construction of the class prototype, and the definition of the
reported factorization discrepancy $0.62$.}

\resp{OpenCLIP \textbf{ViT-B-32}, checkpoint \texttt{laion2b\_s34b\_b79k}. CIFAR-10
\textbf{train} split, $6{,}000$ images drawn with a fixed generator seed. Preprocessing is
OpenCLIP's own transform for that checkpoint; image embeddings are $L_2$-normalised. The source
is the embedding of a Gaussian-blurred image, kernel size $9$, $\sigma$ drawn uniformly from
$[0.8,3.0]$ per call. The condition is the $L_2$-normalised mean of the clean embeddings of the
image's own class. Split $80/20$; the flow is width $512$, three blocks, $\eta=10^{-4}$,
$3{,}000$ steps, batch $256$, sampled with $32$ Heun steps. Target alignment is the mean cosine
between the output embedding and the clean target embedding, evaluated on the held-out $20\%$.}

\interp{Everything except the last item is settled. \textbf{The $0.62$ is the second untraceable
number}: no quantity in the code is named a factorization discrepancy. The nearest candidate is
the estimated off-condition component $\hat\beta=0.6010$, which is close but not equal, and which
already appears in the paper under its own name and role. The value should either be identified
against a named column of the run that produced it, or the sentence removed. The re-run of this
stage has just completed on the cluster and I will check its output before we decide.

One further point that belongs in this section: an earlier version of the label probe was trained
and scored on the same $6{,}000$ embeddings, which with a $512$-dimensional feature and ten
classes is roughly one parameter per sample. That entry was withdrawn as instrumentation rather
than a finding. The current implementation uses a disjoint $80/20$ split and reports the held-out
accuracy alongside the in-sample one.}

\verdict{\NEU}{\FIN}
\end{todoitem}

% ------------------------------------------------------------------
\begin{todoitem}{7}{lines 162--164}
\ask{The 27 empirical threshold estimates $\widehat\beta^\star(A,j)$, or the underlying measured
$\Delta(\beta)$ values over the tested $\beta$-grid.}

\resp{Both are attached. The sweep covers $A\in\{0.15,0.25,0.37,0.50,0.70,0.90\}$ and
$j\in\{0,0.05,0.15,0.30,0.50,0.80\}$, that is $36$ cells, of which $27$ have a threshold defined
under the existence criterion; each cell is measured over
$\beta\in\{0,0.2,0.4,0.6,0.7,0.8,0.9,0.95\}$ and $\widehat\beta^\star$ obtained by linear
interpolation at the sign change of $\Delta(\beta)$. The per-cell table is printed in
\texttt{answers.json} under \texttt{threshold\_grid}; \texttt{beta\_star.csv} carries the raw
$\Delta(\beta)$ in its \texttt{delta} column, so the full comparison can be rebuilt directly.

Maximum absolute error over the $27$ cells: $\mathbf{0.0334}$.}

\interp{This reproduces the $3.3\times10^{-2}$ quoted in the draft to three decimals from the raw
data, independently of the original analysis path. It is one of the two least contested results
in the document.}

\verdict{\POS}{\FIN}
\end{todoitem}

% ------------------------------------------------------------------
\begin{todoitem}{9}{lines 188--191}
\ask{Whether the capacity sweep has now been repeated at a second learning rate.}

\resp{It has been repeated at a \textbf{third}. The grid is now
$\eta\in\{10^{-4},3\times10^{-4},10^{-3}\}$ $\times$ four widths $\times$ five seeds, $120$
trainings, identical protocol. All four registered predictions hold.

\begin{center}
\small
\begin{tabular}{@{}rrrrr@{}}
\toprule
$\eta$ & width $256$ & width $512$ & width $1024$ & width $2048$\\
\midrule
$10^{-4}$          & $0.0419$ & $0.0524$ & $0.0558$ & $0.0612$\\
$3\times10^{-4}$   & $0.0835$ & $0.0938$ & $0.0873$ & $0.0767$\\
$10^{-3}$          & $0.0892$ & $0.1158$ & $0.1020$ & $0.1018$\\
\bottomrule
\end{tabular}
\end{center}

The sign holds in all twelve cells, at a minimum of $16.2$ standard errors. The mean gap grows
monotonically with the learning rate at \emph{every} width: $0.0528$, $0.0853$, $0.1022$,
Spearman $0.90$.}

\interp{This repairs the weakest entry in the register. The previous prediction ``the gap is
learning-rate independent'' was contradicted and stood as a bare falsification; it is now
replaced by a characterised monotone trend, which is what an initial-time amplification predicts
and a generic optimisation failure does not.

The differential is the substantive part and should be stated explicitly. Raising the step size
from $10^{-4}$ to $10^{-3}$ costs $\M_1$ between $1\%$ and $7\%$ of its terminal fidelity and
costs $\M_2$ between $19\%$ and $27\%$: the source-conditioned arm is \textbf{four to sixteen
times more sensitive to step size} than the marginal one, at matched architecture and objective.

One caveat to report rather than round away: at $\eta=3\times10^{-4}$ the gap does narrow with
width, $0.0835\rightarrow0.0767$, which is $1.9$ standard errors --- just inside our two-sigma
bar, so the registered claim passed, but a reader who recomputes will find it. The text should
say ``at no learning rate does the widest model close the gap at the two-standard-error level''.

Finally, the diversity collapse is universal across all twelve cells and \emph{worsens} with
capacity: the source-conditioned pairwise cosine runs $0.344\rightarrow0.506$ from width $256$ to
$2048$ at $\eta=10^{-4}$, against a true-posterior reference of $0.050$ and a marginal-model mean
of $0.225$.}

\verdict{\POS}{\FIN}
\end{todoitem}

% ------------------------------------------------------------------
\begin{todoitem}{10}{lines 224--226}
\ask{The full terminal-fidelity values for each $\varepsilon$, together with the definition and
numerical values of the amplification integral used in the Spearman calculation.}

\resp{\begin{center}
\small
\begin{tabular}{@{}rrrrr@{}}
\toprule
$\varepsilon$ & $\int_\varepsilon^1 t^{-1}dt$ & cosine $\uparrow$ & polar $W_1$ $\downarrow$ &
pairwise cosine $\downarrow$\\
\midrule
$0.00$ & $7.304$ & $0.2166$ & $0.0367$ & $0.6561$\\
$0.01$ & $4.610$ & $0.2227$ & $0.0459$ & $0.6574$\\
$0.02$ & $3.913$ & $0.2286$ & $0.0664$ & $0.6595$\\
$0.05$ & $2.996$ & $0.2528$ & $0.1880$ & $0.6757$\\
$0.10$ & $2.303$ & $0.2827$ & $0.3782$ & $0.7268$\\
$0.20$ & $1.609$ & $0.2997$ & $0.5807$ & $0.8233$\\
$0.35$ & $1.050$ & $0.2945$ & $0.6939$ & $0.9028$\\
\bottomrule
\end{tabular}
\end{center}

Spearman between the integral and the negated cosine: $\mathbf{0.9643}$. The integral diverges at
$\varepsilon=0$; that point is included by flooring the lower limit at $10^{-3}$ and evaluating a
$400$-point trapezoid on a uniform grid, which returns $7.304$ rather than the analytic
$\ln 10^{3}=6.908$ because $1/t$ is steep near the floor. Since Spearman is rank-based, any floor
below the smallest positive tested $\varepsilon$ leaves the ranking unchanged.}

\interp{\textbf{This obliges a change to the draft.} The current text reads ``truncating the
integration recovers most of it, $0.2166\rightarrow0.2997$''. That is true in terminal cosine and
false on both distributional instruments: over the same sweep $W_1$ rises from $0.0367$ to
$0.6939$, a factor of nineteen, and pairwise cosine from $0.656$ to $0.903$. For reference, doing
nothing scores $W_1=0.7667$ and pairwise $0.978$, so truncating to $\varepsilon=0.35$ carries the
model \textbf{ninety per cent of the way back to the do-nothing baseline} while the cosine
improves.

Truncation is therefore a \emph{probe}, not a remedy. It confirms the diagnosis --- the damage is
concentrated near $t=0$, and the recovery in cosine tracks the amplification integral at
$0.964$ --- but the word ``recovers'' has to go. This is the same mean-seeking inversion as
item~2, appearing for the third time, and reporting it strengthens the paper's own thesis rather
than weakening it.}

\verdict{\COR}{\FIN}
\end{todoitem}

% ------------------------------------------------------------------
\begin{todoitem}{11}{lines 227--229}
\ask{The complete auxiliary-source perturbation sweep, all tested values of $\sigma_{\mathrm{aux}}$,
and the exact construction of the tangential perturbation used during training and sampling.}

\resp{The auxiliary input is $a=\exp_{z_0}(\sigma_{\mathrm{aux}}\,u)$ with $u$ a uniformly
distributed unit tangent vector at $z_0$: a random tangent direction, scaled, then carried back to
the sphere by the exponential map. \textbf{The same law is applied during training and during
sampling}, which is the point of the question; the sampler is started at $z_0$ in every case.

\begin{center}
\small
\begin{tabular}{@{}rrrr@{}}
\toprule
$\sigma_{\mathrm{aux}}$ & cosine $\uparrow$ & polar $W_1$ $\downarrow$ & pairwise cosine $\downarrow$\\
\midrule
$0.00$ & $0.2149$ & $0.0367$ & $0.6577$\\
$0.10$ & $0.2435$ & $0.0425$ & $0.3563$\\
$0.30$ & $0.2958$ & $0.0453$ & $0.2438$\\
$0.60$ & $0.3095$ & $0.0404$ & $0.2221$\\
$1.00$ & $0.3105$ & $0.0354$ & $0.2076$\\
\bottomrule
\end{tabular}
\end{center}

One caution for reproducibility: the constant at the head of the notebook lists eight values of
$\sigma_{\mathrm{aux}}$, but the stage that produced this table iterates over the five listed
above. The five are what should be reported.}

\interp{This is the constructive result of the paper and it is currently understated.
Auxiliary noise improves \emph{all three} instruments simultaneously: cosine rises to $0.3105$
against the marginal model's $\approx0.318$, $W_1$ falls to $0.0354$, and pairwise cosine falls
from $0.658$ to $0.208$ against the marginal model's $0.183$. Where truncation buys cosine by
destroying the distribution, noising the auxiliary input recovers the marginal model on every
axis at once.

Together with item~10 the section becomes \emph{one diagnosis, two probes, one remedy}, and the
remedy is a deployable recipe: if the source must be supplied to the field, supply a noised
version of it.}

\verdict{\POS}{\FIN}
\end{todoitem}

% ------------------------------------------------------------------
\begin{todoitem}{15}{lines 319--322}
\ask{The completed super-resolution experiment: the three tested resolutions, the corresponding
theoretical dimensional bounds, the measured source and target dimensions, and the dimensionality
estimation procedure.}

\resp{STL-10, $5{,}000$ training and $1{,}024$ test images, resized and centre-cropped to each
resolution. The source is $z_0=\mathsf{UD}z_1$, average pooling by $f=4$ followed by
nearest-neighbour upsampling; the condition $c$ is the class label. Conditioned on $c$ the source
is a deterministic function of the target, so the endpoints are maximally dependent. Rank is
estimated from the Gram matrix in float64 with strictly more samples than the bound.

\begin{center}
\small
\begin{tabular}{@{}rrrrr@{}}
\toprule
resolution & ambient & bound $r=3(\mathrm{res}/f)^2$ & measured rank $z_0$ & measured rank $z_1$\\
\midrule
$32$ & $3072$  & $192$  & $192$  & $1023$\\
$48$ & $6912$  & $432$  & $432$  & $1983$\\
$96$ & $27648$ & $1728$ & $1728$ & $3071$\\
\bottomrule
\end{tabular}
\end{center}

The measured source rank equals the predicted bound exactly at all three resolutions, against
target ranks that are sample-limited full.

Additionally, the distributional comparison that the draft records as not run has now been
carried out, at two resolutions so far:

\begin{center}
\small
\begin{tabular}{@{}lrrrrrr@{}}
\toprule
& \multicolumn{3}{c}{resolution $32$} & \multicolumn{3}{c}{resolution $48$}\\
Arm & PSNR $\uparrow$ & energy $\downarrow$ & $W_1^{\mathrm{hf}}$ $\downarrow$
    & PSNR $\uparrow$ & energy $\downarrow$ & $W_1^{\mathrm{hf}}$ $\downarrow$\\
\midrule
identity & $19.74$ & $6.743$ & $154.0$ & $20.77$ & $8.942$ & $270.1$\\
$\M_1$   & $21.49$ & $0.734$ & $92.8$  & $22.99$ & $0.826$ & $154.9$\\
$\M_2$   & $21.56$ & $1.017$ & $104.1$ & $23.09$ & $0.959$ & $164.5$\\
\midrule
null floor & --- & $-0.003$ & $7.0$ & --- & $0.001$ & $12.3$\\
\bottomrule
\end{tabular}
\end{center}

Here ``energy'' is the energy distance between the distribution of null-space residuals of the
outputs and of the targets, and $W_1^{\mathrm{hf}}$ the one-dimensional Wasserstein distance
between their high-frequency energy distributions; the null floor is a second independent draw
from the target set.}

\interp{Two separate results. The rank measurement is the strongest single number in the paper:
an exact hit on a predicted bound, three times, with no trained model involved. It establishes
that the source support has strictly smaller dimension than the target support, which is what
forbids a Lipschitz transport between them.

The distributional table settles the model ordering that the draft leaves open. At both
resolutions $\M_2$ attains the higher PSNR while fitting the distribution of restored detail
worse on \emph{both} instruments, consistently across all three seeds. The reversal in the
restoration setting is therefore the mean-seeking inversion again, measured rather than inferred,
and should be reported as a third instance of it and not as support for the ordering.

Scope caution: the claim should be drawn only between $\M_1$ and $\M_2$, which are matched in
capacity and training. The regression arm records the best energy distance of all learned arms
together with the worst PSNR, so the statistic is not a collapse index on its own.}

\verdict{\POS}{\RUN{} --- resolutions $32$ and $48$ complete; $96$ is running now, as is the
six-setting source-informativeness sweep}
\end{todoitem}

% ------------------------------------------------------------------
\section*{Two results not requested in the red items}

\paragraph{The $6.6\times$ ratio is now derived.} It was recorded as measured, not predicted. The
span response carries two suppression factors the skew response does not --- a window $t(1-t)$
vanishing at both endpoints, and the rate at which the two interpolant variances separate. With
homogeneous marginals every symmetric response is exactly zero at first order, while the skew
response has a closed form. The ratio is consequently \emph{not} a universal constant: it
diverges as the marginals homogenise, running from $76.8$ to $7.8$ as we interpolate one marginal
toward the other, so $6.585$ is the sampled ensemble's value at its own separation. This also
explains the registered contradiction ``skew and span equally fragile''. Full note attached
(\texttt{ratio.tex}), with six numerical checks against the same integrator.

\paragraph{The criterion holds on a model we did not train.} Stable Diffusion v1.5 image-to-image,
STL-10 photographs degraded by downsample--upsample, scored by CLIP alignment to the clean image:

\begin{center}
\small
\begin{tabular}{@{}rrr@{}}
\toprule
degradation factor & source alignment & best achievable gain from transporting\\
\midrule
$2$  & $0.997$ & $-0.189$\\
$4$  & $0.853$ & $-0.061$\\
$8$  & $0.733$ & $-0.051$\\
$16$ & $0.430$ & $+0.115$\\
\bottomrule
\end{tabular}
\end{center}

Monotone in degradation, crossing zero between $8\times$ and $16\times$: transport begins to pay
once the source's alignment to the target falls below roughly $0.64$. At $2\times$ no tested
strength beats leaving the image alone, the worst case losing $0.403$. The do-nothing regime is
therefore reachable in a deployed system and not an artefact of the synthetic model. The
non-obvious cases are $4\times$ and $8\times$, where the source is visibly degraded, a
practitioner would run the model without hesitating, and running it still makes the result worse.

This experiment is $64$ images, one checkpoint and one degradation family; it is a targeted test,
not a broad validation, and should be described as such.

\end{document}
